\documentclass[10pt,a4paper]{article}

% Academic CV for faculty applications in Turkey (Bilkent, IYTE).
% Reorganises resume.tex and the _publications/ records; it adds no facts
% beyond the Google Scholar figures the owner supplied on 2026-11-01.
% CI builds it beside resume.tex into Ozgur_Ural_Academic_CV.pdf at the site
% root, linked only from the /publications/ lead (out of the way on purpose:
% the site serves industry readers). Check a change locally with
%   latexmk -pdf -interaction=nonstopmode -halt-on-error -outdir=<dir> academic_cv.tex

\usepackage[a4paper,margin=0.6in]{geometry}
\usepackage{enumitem}
\usepackage{parskip}
\usepackage{mathpazo}
\usepackage{textcomp}
\usepackage[T1]{fontenc}
\usepackage[utf8]{inputenc}
\usepackage[english]{babel}
\usepackage[hidelinks]{hyperref}
\usepackage{needspace}
\usepackage{xcolor}

% Same reasons as resume.tex: no hyphenation, so an ATS or a search finds
% whole words.
\hyphenpenalty=10000
\emergencystretch=2em

\setlist[itemize]{itemsep=3pt, parsep=0pt, topsep=3pt, leftmargin=*}
\setlist[enumerate]{itemsep=4pt, parsep=0pt, topsep=3pt, leftmargin=*, before=\raggedright}
\definecolor{light}{gray}{0.55}
\newcommand{\spacedhrule}[2]{\vspace{#1}\hrule\vspace{#2}}
\newcommand{\cvsection}[1]{\par\needspace{4\baselineskip}\vspace{0.5em}{\large\textbf{#1}}\par\vspace{0.1em}\hrule\vspace{0.35em}}
\newcommand{\cvsub}[1]{\par\needspace{3\baselineskip}\vspace{0.2em}\textit{#1}\par}
\long\def\entry#1#2#3#4{\par\needspace{4\baselineskip}\noindent\textbf{#1}\hfill{#2}\par\nopagebreak\noindent{#3}\par\nopagebreak#4\par\vspace{0.25em}}

\begin{document}
\sloppy
\renewcommand\labelitemi{\textbullet}

\noindent{\LARGE\textbf{Ozgur Ural, Ph.D.}}\hfill{\small\color{light}Curriculum Vitae, \today}\par
\vspace{0.3em}
\noindent Leiden, Netherlands \quad\textbullet\quad \href{mailto:drozgurural@gmail.com}{drozgurural@gmail.com} \quad\textbullet\quad +31 6 2903 0434\par
\noindent \href{https://ozgurural.github.io/}{ozgurural.github.io} \quad\textbullet\quad \href{https://scholar.google.com/citations?user=lYx8fqsAAAAJ&hl=en&sortby=pubdate}{Google Scholar} \quad\textbullet\quad ORCID \href{https://orcid.org/0000-0003-1329-4303}{0000-0003-1329-4303} \quad\textbullet\quad \href{https://linkedin.com/in/uralozgur}{linkedin.com/in/uralozgur}\par

\spacedhrule{0.4em}{0.4em}

Machine learning researcher working on \textbf{trustworthy AI for mission-critical systems}: evidence of how a model was trained and who owns it, robustness against attackers, and predictable timing when a model runs inside a real-time loop. Ph.D. in Electrical Engineering and Computer Science, Embry-Riddle Aeronautical University (2025), with a machine-learning dissertation on Proof-of-Learning and model watermarking and three first-author articles in IEEE Access. More than twelve years of professional experience in industry and research, including production systems in cybersecurity, defense and flight simulation; currently at Avion Full Flight Simulators, leading the company's enterprise AI strategy.

\cvsection{RESEARCH INTERESTS}
\begin{itemize}
  \item Trustworthy and verifiable machine learning: Proof-of-Learning, model watermarking, model provenance and ownership verification
  \item Security of machine learning: spoofing attacks on training records, adversarial robustness
  \item Blockchain-enhanced machine learning
  \item Machine learning in real-time and safety-critical systems
  \item Verifiable oversight of autonomous software agents
  \item Natural language processing for cyber security
\end{itemize}

\cvsection{EDUCATION}

\entry{Ph.D. in Electrical Engineering and Computer Science}{Aug 2021 -- Aug 2025}{Embry-Riddle Aeronautical University, Daytona Beach, Florida, USA}{%
\begin{itemize}
  \item Dissertation: \href{https://commons.erau.edu/edt/905/}{Enhancing Proof-of-Learning Security Against Spoofing Attacks Using Model Watermarking}.
  \item Advisor: Dr. Kenji Yoshigoe (Committee Chair); committee included IEEE Fellow Dr. Houbing Song.
\end{itemize}}

\entry{M.Sc. in Cyber Security}{2015 -- 2019}{Middle East Technical University (METU), Ankara, Turkey}{%
\begin{itemize}
  \item Thesis: \href{https://open.metu.edu.tr/handle/11511/43747}{Automatic Detection of Cyber Security Events from Turkish Twitter Stream and Turkish Newspaper Data}.
  \item Advisor: Prof. Cengiz Acart\"urk.
\end{itemize}}

\entry{B.Sc. in Computer Engineering}{2010 -- 2014}{Middle East Technical University (METU), Ankara, Turkey}{%
\begin{itemize}
  \item English-medium programme; entered among roughly the top 1,000 of more than 2 million national candidates.
\end{itemize}}

\entry{Bachelor of Business Administration}{2012 -- 2017}{Anadolu University, Eskisehir, Turkey}{}

\cvsection{PUBLICATIONS}

\noindent \href{https://scholar.google.com/citations?user=lYx8fqsAAAAJ&hl=en&sortby=pubdate}{Google Scholar} (November 2026): 53 citations, h-index 4.

\cvsub{Journal Articles}
\begin{enumerate}
  \item[{[J3]}] \textbf{Ural, O.} and Yoshigoe, K. (2025). \href{https://doi.org/10.1109/ACCESS.2025.3642198}{SecurePoL: Integration of Watermarking With Proof-of-Learning to Enhance Security Against Spoofing Attacks}. \textit{IEEE Access}, vol. 13, pp. 213067--213091. DOI: 10.1109/ACCESS.2025.3642198. \textit{Research article.}
  \item[{[J2]}] \textbf{Ural, O.} and Yoshigoe, K. (2024). \href{https://doi.org/10.1109/ACCESS.2024.3489776}{Enhancing Security of Proof-of-Learning Against Spoofing Attacks Using Feature-Based Model Watermarking}. \textit{IEEE Access}, vol. 12, pp. 169567--169591. DOI: 10.1109/ACCESS.2024.3489776. \textit{Research article.}
  \item[{[J1]}] \textbf{Ural, O.} and Yoshigoe, K. (2023). \href{https://doi.org/10.1109/ACCESS.2023.3344669}{Survey on Blockchain-Enhanced Machine Learning}. \textit{IEEE Access}, vol. 11, pp. 145331--145362. DOI: 10.1109/ACCESS.2023.3344669. \textit{Survey article.}
\end{enumerate}
\noindent{\small IEEE Access: 2025 Journal Impact Factor 4.2 (JCR 2026).}

\cvsub{Conference Papers}
\begin{enumerate}
  \item[{[C2]}] \textbf{Ural, O.} and Acart\"urk, C. (2021). \href{https://doi.org/10.5220/0010201600660076}{Automatic Detection of Cyber Security Events from Turkish Twitter Stream and Newspaper Data}. In \textit{Proceedings of the 7th International Conference on Information Systems Security and Privacy (ICISSP 2021)}, pp. 66--76. DOI: 10.5220/0010201600660076.
  \item[{[C1]}] Yama\c{c}, A. \"O., Do\c{c}i, J., \textbf{Ural, O.}, \"Ozarslan, \"O. F. and Temizer, S. (2014). \href{https://senior.ceng.metu.edu.tr/2014/clover/}{Clover: Autonomous Cargo and Mail Delivery}. \textit{Turkish Autonomous Robots Conference (TORK 2014)}, Ankara, Turkey.
\end{enumerate}

\cvsub{Theses}
\begin{enumerate}
  \item[{[T2]}] \textbf{Ural, O.} (2025). \href{https://commons.erau.edu/edt/905/}{Enhancing Proof-of-Learning Security Against Spoofing Attacks Using Model Watermarking}. Doctoral dissertation, Embry-Riddle Aeronautical University.
  \item[{[T1]}] \textbf{Ural, O.} (2019). \href{https://open.metu.edu.tr/handle/11511/43747}{Automatic Detection of Cyber Security Events from Turkish Twitter Stream and Turkish Newspaper Data}. Master's thesis, Middle East Technical University.
\end{enumerate}

\cvsub{Technical Reports (not peer reviewed)}
\begin{enumerate}
  \item[{[R1]}] \textbf{Ural, O.} and Erdur, E. (2016). \href{https://doi.org/10.13140/RG.2.2.24058.08649}{Secure Proxy on Cloud}. Technical report. DOI: 10.13140/RG.2.2.24058.08649.
\end{enumerate}

\noindent Narrated explainers of this research: \href{https://ozgurural.github.io/lab/}{ozgurural.github.io/lab}.

\cvsection{RESEARCH EXPERIENCE}

\entry{Graduate Research Assistant}{Sep 2021 -- Dec 2024, May 2025 -- Oct 2025}{Embry-Riddle Aeronautical University, Daytona Beach, Florida, USA. Advisor: Dr. Kenji Yoshigoe}{%
\begin{itemize}
  \item Original research on verifiable, trustworthy machine learning across Proof-of-Learning, model watermarking and adversarial robustness, published as three first-author IEEE Access articles [J1--J3] and the doctoral dissertation [T2].
\end{itemize}}

\entry{M.Sc. Thesis Research}{2015 -- 2019}{Middle East Technical University, Ankara, Turkey. Advisor: Prof. Cengiz Acart\"urk}{%
\begin{itemize}
  \item Automatic detection of cyber security events from Turkish Twitter and newspaper data [T1], published at ICISSP 2021 [C2].
\end{itemize}}

\cvsection{PROFESSIONAL SERVICE}
\begin{itemize}
  \item \textbf{Program Committee (2026):} NLPAICS 2026, 2nd Workshop on NLP Applied to Information and Cyber Security, Alicante.
  \item \textbf{Reviewer:} IEEE Access and other journals; IEEE and ACM conferences (security, machine learning, distributed systems).
\end{itemize}

\cvsection{AWARDS}
\begin{itemize}
  \item First Prize, TUBITAK national project competition (2014).
  \item Second Prize, METU Computer Engineering graduation Demo Day (2014).
\end{itemize}

\cvsection{INDUSTRY EXPERIENCE}

\entry{Senior Software Engineer}{Oct 2023 -- Present}{Avion Full Flight Simulators, Nieuw-Vennep, Netherlands}{%
\begin{itemize}
  \item Lead Avion's enterprise AI strategy: authored the AI adoption roadmap and a three-tier architecture spanning air-gapped on-prem LLMs, a governed cloud-LLM API tier, and edge deployment.
  \item Won approval for two R\&D projects under WBSO, the Dutch government's R\&D tax-credit programme assessed by the Netherlands Enterprise Agency: sole person responsible for preparing and architecting both, for the Avion Cloud Environment and for autonomous engineering agents with verifiable oversight.
  \item Built Avion Assistant, an engineering-knowledge copilot offered to staff as a proof of concept, with hybrid retrieval, every answer cited to its source page, and regression checks on each release; added a prototype agent foundation that records its actions.
  \item Designed the strategy's proposed AI workflows, each keeping a person accountable for the outcome; built the strategy portal end to end.
  \item Principal developer of the Avion Cloud Environment (about 70\% of commits since 2024), a platform to configure, monitor and diagnose simulator bus systems (SvelteKit, TypeScript, Python, gRPC, Redis, PostgreSQL; Scala data senders).
  \item Built the Avion Debriefing System: records a full-flight-simulator training session (displays, cockpit camera, instructor station, audio), then packages, encrypts and transfers it (C\#/.NET 8, PowerShell, MSI installer); installed and verified on a simulator, with user and acceptance-test manuals.
  \item Built NavDB Sync, version control for simulator navigation databases: operators keep their custom data through 28-day AIRAC updates, and merges run on the server so a conflict never reaches a simulator (Python, GitLab CI/CD).
  \item Developed real-time simulator software in C, C++ and Scala, including a low-latency telemetry pipeline with batching and bounded-queue backpressure.
\end{itemize}}

\entry{Software Team Lead}{Nov 2020 -- Aug 2021}{Havelsan, Ankara, Turkey}{%
\begin{itemize}
  \item Led a team of 14 engineers (developers, QA, DevOps, support) delivering the Havelsan DLP endpoint Data Leakage Prevention product; developed core endpoint-agent features in C, C++, Java and Go.
\end{itemize}}

\entry{Expert Software Engineer}{Feb 2019 -- Nov 2020}{STM Defence Technologies, Ankara, Turkey}{%
\begin{itemize}
  \item Developed mission-control and ground-control components for the Kargu and Togan autonomous UAVs in C++11 with Boost and Qt, under hard-real-time, safety-critical constraints; designed Power-On and Continuous Built-In Tests.
\end{itemize}}

\entry{Expert Software Engineer}{Jul 2014 -- Feb 2019}{Comodo Cybersecurity, Ankara, Turkey}{%
\begin{itemize}
  \item Led design and development of Windows services and applications in C++17 with Boost and WinAPI (Comodo Dome Secure Web Gateway, Comodo Patch Manager, Comodo Dragon browser); established unit testing, code review and Jenkins CI/CD.
\end{itemize}}

\cvsection{TECHNICAL SKILLS}
\begin{itemize}
  \item \textbf{Machine learning:} PyTorch, TensorFlow, scikit-learn, deep learning, LLMs, retrieval-augmented generation, federated learning; model watermarking, Proof-of-Learning, adversarial robustness
  \item \textbf{Languages:} C, C++, C\#, Python, TypeScript, JavaScript, Go, Java, Scala
  \item \textbf{Systems:} real-time and distributed systems, multithreading, low-latency pipelines, Linux, gRPC, Protobuf, Redis, PostgreSQL, Docker, GitLab CI/CD, Jenkins, Git
\end{itemize}

\end{document}
